\section{2012 Analysis \& Differential Equations}

\subsection{Individual}

\begin{exercise}
\label{analysis:2012:1}
Compute the integral
\[
\int_0^{\infty} \frac{x^p}{1+x^2}\,dx, \quad -1<p<1.
\]
\end{exercise}

\begin{solution}
\textbf{Prerequisites:} 
\begin{itemize}
    \item \textbf{Complex Analysis:} Cauchy's Residue Theorem, which states that the integral around a closed contour is $2\pi i$ times the sum of residues at poles inside.
    \item \textbf{Branch Cuts:} Understanding that $x^p = e^{p \ln x}$ is multi-valued for non-integer $p$, requiring a branch cut (usually $[0, \infty)$) to make it a single-valued function.
    \item \textbf{Improper Integrals:} Convergence criteria for integrals involving power functions.
\end{itemize}

\bigskip

\textbf{Steps:}
\begin{enumerate}
    \item \textbf{Define the Complex Function:} Define $f(z) = \frac{z^p}{1+z^2}$. To handle $z^p$, we choose a branch of the logarithm such that $0 < \arg z < 2\pi$. This places a branch cut along the positive real axis.
    
    \item \textbf{Set up the Contour:} We use a "keyhole contour" $C$ consisting of:
    \begin{itemize}
        \item $L_1$: A segment just above the real axis from $\epsilon$ to $R$.
        \item $C_R$: A large circle of radius $R$ counter-clockwise.
        \item $L_2$: A segment just below the real axis from $R$ to $\epsilon$.
        \item $C_\epsilon$: A small circle of radius $\epsilon$ clockwise around the origin.
    \end{itemize}

    \item \textbf{Analyze the Circles:} 
    \begin{itemize}
        \item On $C_R$, $|f(z)| \approx R^p/R^2 = R^{p-2}$. Since $p < 1$, the integral $\int_{C_R} f(z)dz \to 0$ as $R \to \infty$.
        \item On $C_\epsilon$, $|f(z)| \approx \epsilon^p$. Since $p > -1$, the integral $\int_{C_\epsilon} f(z)dz \to 0$ as $\epsilon \to 0$.
    \end{itemize}

    \item \textbf{Analyze the Segments:} 
    \begin{itemize}
        \item On $L_1$, $z = x e^{i0}$, so $f(z) = \frac{x^p}{1+x^2}$. The integral is $I$.
        \item On $L_2$, $z = x e^{i2\pi}$, so $f(z) = \frac{x^p e^{i2\pi p}}{1+x^2}$. The integral is $\int_R^\epsilon \frac{x^p e^{i2\pi p}}{1+x^2} dx \to -e^{i2\pi p} I$.
    \end{itemize}

    \item \textbf{Calculate Residues:} The poles of $1+z^2=0$ are at $z = i = e^{i\pi/2}$ and $z = -i = e^{i3\pi/2}$.
    \begin{itemize}
        \item $\text{Res}(f, e^{i\pi/2}) = \frac{(e^{i\pi/2})^p}{2e^{i\pi/2}} = \frac{e^{ip\pi/2}}{2i}$.
        \item $\text{Res}(f, e^{i3\pi/2}) = \frac{(e^{i3\pi/2})^p}{2e^{i3\pi/2}} = \frac{e^{i3p\pi/2}}{-2i}$.
    \end{itemize}

    \item \textbf{Sum and Solve:} By the Residue Theorem:
    \[ I(1 - e^{i2\pi p}) = 2\pi i \left( \frac{e^{ip\pi/2} - e^{i3p\pi/2}}{2i} \right) = \pi (e^{ip\pi/2} - e^{i3p\pi/2}). \]
    Divide by $(1 - e^{i2\pi p})$ and simplify using Euler's formula to get:
    \[ I = \frac{\pi}{2\cos(p\pi/2)}. \]
\end{enumerate}

\bigskip

\textbf{Intuition:} 
The function $x^p$ is tricky because it's not a standard polynomial; it "wraps" around the complex plane. By using a keyhole contour, we catch the "jump" the function makes when you go all the way around the origin ($2\pi$ radians). This jump, combined with the residues at the vertical poles $\pm i$, allows us to extract the real integral without actually finding an antiderivative.
\end{solution}

\begin{exercise}
\label{analysis:2012:2}
Construct a one to one conformal mapping from the region
\[
U = \{z\in\mathbb{C} \mid |z-\frac{i}{2}|<\frac{1}{2}\} \setminus \{z \mid |z-\frac{i}{4}|<\frac{1}{4}\}
\]
onto the unit disk.
\end{exercise}

\begin{solution}
\textbf{Prerequisites:}
\begin{itemize}
    \item \textbf{M\"obius Transformations:} Transformations of the form $\frac{az+b}{cz+d}$ that map circles/lines to circles/lines.
    \item \textbf{Inversion:} The map $z \mapsto 1/z$, which is a special M\"obius transform that maps points near the origin to points near infinity.
    \item \textbf{Standard Regions:} Knowing how to map strips to half-planes (exponential) and half-planes to disks (M\"obius).
\end{itemize}

\bigskip

\textbf{Steps:}
\begin{enumerate}
    \item \textbf{Analyze the Geometry:} The region $U$ is the space between two circles that both touch the origin $(0,0)$. The outer circle has radius $1/2$ (centered at $i/2$) and the inner has radius $1/4$ (centered at $i/4$).
    
    \item \textbf{Send Tangency to Infinity:} Use $w_1 = 1/z$. Since both circles pass through the origin, they will map to parallel lines in the $w_1$-plane. 
    Specifically, the outer circle $x^2 + (y-1/2)^2 = 1/4 \implies x^2+y^2-y=0$ becomes $\text{Im}(w_1) = -1$.
    The inner circle $x^2+(y-1/4)^2=1/16 \implies x^2+y^2-y/2=0$ becomes $\text{Im}(w_1) = -2$.
    The region $U$ maps to the horizontal strip $-2 < \text{Im}(w_1) < -1$.

    \item \textbf{Normalize the Strip:} We want to map this to a strip centered at the origin of width $\pi$.
    Let $w_2 = \pi(w_1 + 1.5i)$. This shifts the strip to $-\pi/2 < \text{Im}(w_2) < \pi/2$.

    \item \textbf{Map to Half-Plane:} Use $w_3 = e^{w_2}$. The exponential map takes a horizontal strip of width $\pi$ to a half-plane. For this specific range, it maps to the right half-plane $\text{Re}(w_3) > 0$.

    \item \textbf{Map to Unit Disk:} Use the standard M\"obius map $w = \frac{w_3 - 1}{w_3 + 1}$ to map the right half-plane to the unit disk $\mathbb{D} = \{w : |w| < 1\}$.
\end{enumerate}

\bigskip

\textbf{Intuition:} 
The "crescent" shape is hard to work with directly. By sending the "pinch point" (the origin) to infinity, the curved boundaries "straighten out" into parallel lines. Once you have a strip, you can use the exponential function—which is periodic—to "unroll" that strip into a full half-plane, which is then easily squashed into a disk.
\end{solution}

\begin{exercise}
\label{analysis:2012:3}
Let $f(x)$ be a nonlinear $C^2$ function on $\mathbb{R}$. Show that
\[
\sup|f'(x)|^2 \leq 4\sup|f(x)|\sup|f''(x)|.
\]
\end{exercise}

\begin{solution}
\textbf{Prerequisites:}
\begin{itemize}
    \item \textbf{Taylor's Theorem:} Specifically the version $f(x+h) = f(x) + hf'(x) + \frac{h^2}{2}f''(\xi)$.
    \item \textbf{Supremum (Sup):} The least upper bound of a set of values (roughly the "maximum" over the whole domain).
    \item \textbf{Optimization:} Finding the minimum of a function of one variable using derivatives.
\end{itemize}

\bigskip

\textbf{Steps:}
\begin{enumerate}
    \item \textbf{Notation:} Let $M_0 = \sup|f|$, $M_1 = \sup|f'|$, and $M_2 = \sup|f''|$. We want to show $M_1^2 \leq 4M_0 M_2$.
    
    \item \textbf{Taylor Expansion:} Write the expansion for $f(x+h)$ where $h$ is any positive number:
    \[ f(x+h) = f(x) + hf'(x) + \frac{h^2}{2} f''(\xi) \]
    Rearrange this to isolate the first derivative:
    \[ hf'(x) = f(x+h) - f(x) - \frac{h^2}{2} f''(\xi) \]

    \item \textbf{Apply Absolute Values and Bounds:}
    \[ h|f'(x)| \leq |f(x+h)| + |f(x)| + \frac{h^2}{2} |f''(\xi)| \]
    Using our $M$ notation, this becomes:
    \[ h|f'(x)| \leq M_0 + M_0 + \frac{h^2}{2} M_2 = 2M_0 + \frac{h^2}{2} M_2 \]
    Dividing by $h$:
    \[ |f'(x)| \leq \frac{2M_0}{h} + \frac{h M_2}{2} \]

    \item \textbf{Optimize $h$:} The right side is a function of $h$. To find the "tightest" bound, we minimize $g(h) = \frac{2M_0}{h} + \frac{h M_2}{2}$.
    Setting $g'(h) = -\frac{2M_0}{h^2} + \frac{M_2}{2} = 0$ gives $h = 2\sqrt{M_0/M_2}$.
    
    \item \textbf{Substitute back:} Plugging this $h$ into the inequality:
    \[ |f'(x)| \leq \frac{2M_0}{2\sqrt{M_0/M_2}} + \frac{2\sqrt{M_0/M_2} M_2}{2} = \sqrt{M_0 M_2} + \sqrt{M_0 M_2} = 2\sqrt{M_0 M_2} \]
    Squaring both sides gives $M_1^2 \leq 4 M_0 M_2$.
\end{enumerate}

\bigskip

\textbf{Intuition:} 
Imagine a car's position $f(x)$. $f'$ is the speed and $f''$ is the acceleration. If the car's position is always between $-M_0$ and $M_0$ (it stays in a small garage), but it's moving very fast ($M_1$ is large), it must decelerate or accelerate incredibly hard ($M_2$ must be large) to avoid hitting the walls. This inequality proves that speed, position range, and acceleration are fundamentally linked.
\end{solution}

\begin{exercise}
\label{analysis:2012:4}
Let $f(x)$ be a real measurable function defined on $[a,b]$. Let $n(y)$ be the number of solutions of the equation $f(x)=y$. Prove that $n(y)$ is a measurable function on $\mathbb{R}$.
\end{exercise}

\begin{solution}
\textbf{Prerequisites:}
\begin{itemize}
    \item \textbf{Measurable Sets:} Understanding Lebesgue measurability in $\mathbb{R}$ and product spaces $\mathbb{R}^n$.
    \item \textbf{Graph of a Function:} The set $G = \{(x, f(x)) \mid x \in D\}$. If $f$ is measurable, $G$ is a measurable subset of the product space.
    \item \textbf{Measurable Projection Theorem:} A deep result stating that the projection of a measurable set (specifically a Borel or Suslin set) onto a lower-dimensional space is Lebesgue measurable.
\end{itemize}

\bigskip

\textbf{Steps:}
\begin{enumerate}
    \item \textbf{Define the Target Sets:} We want to show $n(y)$ is measurable. A common technique is to show that for every $k \in \{1, 2, \dots\}$, the set $E_k = \{y \in \mathbb{R} \mid n(y) \geq k\}$ is measurable. This set contains all $y$-values that are "hit" by $f$ at least $k$ times.
    
    \item \textbf{Construct the Configuration Space:} Consider the set $S_k \subset [a, b]^k \times \mathbb{R}$ defined by:
    \[ S_k = \{(x_1, \dots, x_k, y) \mid f(x_i) = y \text{ for } i=1,\dots,k \text{ and } x_i \neq x_j \text{ for } i \neq j \} \]
    This set represents all possible "collections of $k$ distinct inputs" that map to the same output $y$.

    \item \textbf{Show $S_k$ is Measurable:} Since $f$ is measurable, the condition $f(x_i) = y$ is a measurable condition in the product space $[a,b]^k \times \mathbb{R}$. The condition $x_i \neq x_j$ is also measurable (it's the complement of the diagonals $x_i = x_j$). Thus, $S_k$ is a measurable set.

    \item \textbf{Project onto the $y$-axis:} The set $E_k$ is exactly the projection of $S_k$ onto the last coordinate (the $y$ coordinate). 
    \[ E_k = \text{Proj}_y(S_k) \]
    By the Measurable Projection Theorem, $E_k$ is a measurable subset of $\mathbb{R}$.

    \item \textbf{Reconstruct the Function:} The function $n(y)$ (the total count) can be written as the sum of the characteristic functions of these sets:
    \[ n(y) = \sum_{k=1}^\infty \chi_{E_k}(y) \]
    Since each $\chi_{E_k}$ is a measurable function and the sum of measurable functions is measurable, $n(y)$ is measurable.
\end{enumerate}

\bigskip

\textbf{Intuition:} 
Imagine the graph of $f$. We are counting how many times a horizontal line at height $y$ crosses the graph. This "cross-section" counting is a bit like the Fubini theorem in reverse. By showing that the set of $y$'s with "at least $k$ crossings" is measurable for every $k$, we essentially "build" the function $n(y)$ out of measurable building blocks.
\end{solution}

\begin{exercise}
\label{analysis:2012:5}
For $1<p,q<\infty$, $\frac{1}{p}+\frac{1}{q}=1$, let $g$ in $L^q$. Consider the linear functional $F$ on $L^p$ given by: $F(f) = \int fg$. Show that $\|F\| = \|g\|_q$.
\end{exercise}

\begin{solution}
\textbf{Prerequisites:}
\begin{itemize}
    \item \textbf{$L^p$ Spaces and Norms:} $\|f\|_p = (\int |f|^p)^{1/p}$.
    \item \textbf{H\"older's Inequality:} $\int |fg| \leq \|f\|_p \|g\|_q$.
    \item \textbf{Dual Norm:} The norm of a functional $F$ is defined as $\|F\| = \sup_{\|f\|_p=1} |F(f)|$.
\end{itemize}

\bigskip

\textbf{Steps:}
\begin{enumerate}
    \item \textbf{Establish the Upper Bound:} By H\"older's inequality:
    \[ |F(f)| = \left| \int fg \right| \leq \int |fg| \leq \|f\|_p \|g\|_q \]
    If we take the supremum over all $f$ with $\|f\|_p=1$, we get:
    \[ \|F\| \leq \|g\|_q \]

    \item \textbf{Construct the "Optimizing" Function:} We want to find a specific $f$ that makes the inequality an equality. We need $fg$ to be positive and $|f|^p$ to be related to $|g|^q$.
    Define $f(x) = \text{sgn}(g(x)) |g(x)|^{q-1}$. 
    
    \item \textbf{Calculate the Norm of $f$:}
    \[ |f|^p = |g|^{p(q-1)} \]
    Since $\frac{1}{p} + \frac{1}{q} = 1$, we have $p(q-1) = q$. So $|f|^p = |g|^q$.
    The $L^p$ norm is:
    \[ \|f\|_p = \left( \int |g|^q \right)^{1/p} = (\|g\|_q^q)^{1/p} = \|g\|_q^{q/p} \]

    \item \textbf{Evaluate the Functional:}
    \[ F(f) = \int f g = \int \text{sgn}(g) |g|^{q-1} g = \int |g|^q = \|g\|_q^q \]
    
    \item \textbf{Final Comparison:} To get a unit vector, use $\hat{f} = f / \|f\|_p$.
    \[ F(\hat{f}) = \frac{\|g\|_q^q}{\|g\|_q^{q/p}} = \|g\|_q^{q - q/p} \]
    Since $q - q/p = q(1-1/p) = q(1/q) = 1$, we have $F(\hat{f}) = \|g\|_q$.
    This proves $\|F\| \geq \|g\|_q$. Combined with step 1, we get $\|F\| = \|g\|_q$.
\end{enumerate}

\bigskip

\textbf{Intuition:} 
Think of the integral $\int fg$ as a "dot product" in an infinite-dimensional space. The norm $\|F\|$ is the length of the vector $g$. H\"older's inequality is like the Cauchy-Schwarz inequality—it says the dot product is maximized when $f$ "points" in the same direction as $g$. Our construction of $f$ is exactly finding that "parallel" vector in the $L^p$ geometry.
\end{solution}

\begin{exercise}
\label{analysis:2012:6}
Show that the Poisson integral formula $u(x) = \frac{2x_n}{n\alpha_n}\int_{\partial\mathbb{R}_+^n} \frac{g(y)}{|x-y|^n}\,dy$ solves the Dirichlet problem in the half-space.
\end{exercise}

\begin{solution}
\textbf{Prerequisites:}
\begin{itemize}
    \item \textbf{Harmonic Functions:} Functions $u$ that satisfy $\Delta u = 0$.
    \item \textbf{Dirichlet Problem:} Finding a harmonic function in a region that matches given values $g$ on the boundary.
    \item \textbf{Identity Approximations:} A sequence of functions that "become" the Dirac Delta function in the limit.
\end{itemize}

\bigskip

\textbf{Steps:}
\begin{enumerate}
    \item \textbf{Define the Kernel:} Let $K(x, y) = \frac{2x_n}{\omega_n |x-y|^n}$ where $\omega_n$ is the surface area of the unit sphere. We must show $u(x) = \int K(x,y)g(y)dy$ is harmonic.
    
    \item \textbf{Verify Harmonicity:} The Laplacian $\Delta_x$ acts only on the $x$ variable. Since $K(x,y)$ is derived from the Green's function (specifically its normal derivative), and the Green's function is harmonic away from its pole, $K(x,y)$ is harmonic for $x_n > 0$. By differentiating under the integral sign, $\Delta u = 0$.

    \item \textbf{Check the Integral Value:} It is a standard result (via a change of variables to polar coordinates) that for any fixed $x$ with $x_n > 0$:
    \[ \int_{\mathbb{R}^{n-1}} K(x, y) dy = 1 \]
    This means $u(x)$ is a "weighted average" of the boundary values $g(y)$.

    \item \textbf{Analyze the Boundary Limit:} As $x$ approaches a point $x_0$ on the boundary $(x_n \to 0)$, the kernel $K(x,y)$ becomes extremely large when $y$ is near $x_0$ and almost zero elsewhere. 
    Because the total integral is always 1, the kernel "picks out" the value $g(x_0)$.
    \[ \lim_{x \to x_0} u(x) = \lim_{x \to x_0} \int K(x,y) g(y) dy = g(x_0) \]
\end{enumerate}

\bigskip

\textbf{Intuition:} 
The Poisson formula is like a "blurring" operation. The values on the boundary are spread out into the interior. The closer you are to the boundary, the less "blurring" happens for the point directly below you, until finally, at the boundary, you see the exact value. Mathematically, the kernel $K$ acts as a "window" that narrows down to a single point as you touch the boundary.
\end{solution}

\subsection{Team}

\begin{exercise}
\label{analysis:2012:7}
Prove that $f(x) = \exp(-\frac{1}{2}x^T A x)$ is in $L^1(\mathbb{R}^n)$ iff $A$ is positive definite. Compute the integral with a linear term $b^T x$.
\end{exercise}

\begin{solution}
\textbf{Prerequisites:}
\begin{itemize}
    \item \textbf{Linear Algebra:} Quadratic forms $x^T A x$ and matrix diagonalization.
    \item \textbf{Multivariable Calculus:} Change of variables in multiple integrals and Jacobian determinants.
    \item \textbf{Gaussian Integrals:} The basic 1D result $\int_{-\infty}^\infty e^{-ax^2} dx = \sqrt{\pi/a}$ for $a > 0$.
\end{itemize}

\bigskip

\textbf{Steps:}
\begin{enumerate}
    \item \textbf{Integrability Analysis:} The function $f(x) = \exp(-\frac{1}{2}x^T A x)$ is in $L^1(\mathbb{R}^n)$ if the integral $\int_{\mathbb{R}^n} f(x) dx$ is finite.
    Since $A$ is symmetric, we can diagonalize it: $A = Q D Q^T$, where $D = \text{diag}(\lambda_1, \dots, \lambda_n)$ and $Q$ is orthogonal.
    Let $x = Qy$, then $dx = dy$ (since $|\det Q| = 1$). The exponent becomes:
    \[ x^T A x = (Qy)^T (QDQ^T) (Qy) = y^T D y = \sum_{i=1}^n \lambda_i y_i^2 \]
    The integral factors:
    \[ \int_{\mathbb{R}^n} e^{-\frac{1}{2}\sum \lambda_i y_i^2} dy = \prod_{i=1}^n \int_{-\infty}^\infty e^{-\frac{1}{2}\lambda_i y_i^2} dy_i \]
    Each 1D integral converges if and only if $\lambda_i > 0$ for all $i$. This is exactly the definition of $A$ being \textbf{positive definite}.

    \item \textbf{Handling the Linear Term:} Consider the exponent $-\frac{1}{2}x^T A x + b^T x$. We "complete the square" in vector form:
    \[ -\frac{1}{2}(x - A^{-1}b)^T A (x - A^{-1}b) + \frac{1}{2}b^T A^{-1}b \]
    Substituting this into the integral:
    \[ I = \int e^{-\frac{1}{2}(x - A^{-1}b)^T A (x - A^{-1}b)} e^{\frac{1}{2}b^T A^{-1}b} dx \]
    The term $e^{\frac{1}{2}b^T A^{-1}b}$ is constant and can be pulled out. Let $z = x - A^{-1}b$:
    \[ I = e^{\frac{1}{2}b^T A^{-1}b} \int_{\mathbb{R}^n} e^{-\frac{1}{2}z^T A z} dz \]

    \item \textbf{Final Computation:} Using the diagonalization result from Step 1, each 1D integral is $\sqrt{2\pi/\lambda_i}$. Multiplying them:
    \[ \prod \sqrt{\frac{2\pi}{\lambda_i}} = \frac{(2\pi)^{n/2}}{\sqrt{\prod \lambda_i}} = \frac{(2\pi)^{n/2}}{\sqrt{\det A}} \]
    Thus, $I = e^{\frac{1}{2}b^T A^{-1}b} \frac{(2\pi)^{n/2}}{\sqrt{\det A}}$.
\end{enumerate}

\bigskip

\textbf{Intuition:} 
A Gaussian function is a bell curve. In one dimension, if the "width" factor is negative, the curve flips upside down and goes to infinity, so it's not integrable. In $n$ dimensions, positive definiteness ensures the curve "decays" in every possible direction. The linear term $b^T x$ just shifts the center of the bell curve away from the origin.
\end{solution}

\begin{exercise}
\label{analysis:2012:8}
In a simply connected region $V \neq \mathbb{C}$, if two biholomorphisms $\phi_1, \phi_2$ agree at two distinct points, show they are identical.
\end{exercise}

\begin{solution}
\textbf{Prerequisites:}
\begin{itemize}
    \item \textbf{Riemann Mapping Theorem:} Every simply connected region $V \subsetneq \mathbb{C}$ is biholomorphic to the unit disk $\mathbb{D}$.
    \item \textbf{Schwarz Lemma:} If $f: \mathbb{D} \to \mathbb{D}$ is holomorphic with $f(0)=0$, then $|f(z)| \leq |z|$. If $|f(z)| = |z|$ for some $z \neq 0$, then $f(z) = e^{i\theta}z$.
    \item \textbf{Fixed Points:} An automorphism of the disk with two fixed points must be the identity.
\end{itemize}

\bigskip

\textbf{Steps:}
\begin{enumerate}
    \item \textbf{Reduce to the Disk:} Let $h = \phi_2^{-1} \circ \phi_1$. This is an automorphism of the region $V$. We are given that $h(a) = a$ and $h(b) = b$ for two distinct points $a, b \in V$.
    
    \item \textbf{Map to the Standard Region:} By the Riemann Mapping Theorem, there exists a biholomorphism $g: V \to \mathbb{D}$. Let $f = g \circ h \circ g^{-1}$. This $f$ is an automorphism of the unit disk $\mathbb{D}$.
    
    \item \textbf{Identify Fixed Points in the Disk:} Let $a' = g(a)$ and $b' = g(b)$. Since $g$ is a bijection and $a,b$ are distinct, $a', b'$ are distinct points in $\mathbb{D}$. Also:
    \[ f(a') = g(h(g^{-1}(g(a)))) = g(h(a)) = g(a) = a' \]
    Similarly, $f(b') = b'$. So $f$ fixes two distinct points in the unit disk.

    \item \textbf{Apply Automorphism Properties:} Any automorphism of $\mathbb{D}$ is a M\"obius transformation. If we perform another M\"obius transform to move $a'$ to $0$, the new map still fixes $0$ and some non-zero point. By the Schwarz Lemma, a map fixing $0$ must be a rotation $z \mapsto e^{i\theta}z$. If it also fixes a non-zero point, then $e^{i\theta} = 1$.
    
    \item \textbf{Conclusion:} Thus $f$ is the identity map on $\mathbb{D}$. This implies $h$ is the identity map on $V$, so $\phi_1 = \phi_2$.
\end{enumerate}

\bigskip

\textbf{Intuition:} 
Biholomorphisms (conformal maps) are very "rigid." They preserve angles and local structure so strictly that they don't have much freedom. If you know where two points go, the "tension" required to keep the map conformal forces every other point into a specific, unique position.
\end{solution}

\begin{exercise}
\label{analysis:2012:9}
Calculate the measure of the set $E \subset [0,1]$ of numbers with no '4' in their decimal expansion.
\end{exercise}

\begin{solution}
\textbf{Prerequisites:}
\begin{itemize}
    \item \textbf{Lebesgue Measure:} The "length" of a set of points.
    \item \textbf{Geometric Series:} The sum $\sum r^k = \frac{1}{1-r}$ for $|r| < 1$.
    \item \textbf{Complements:} It's often easier to measure the set of things we \textit{don't} want and subtract from the total.
\end{itemize}

\bigskip

\textbf{Steps:}
\begin{enumerate}
    \item \textbf{Identify the Complement:} Let $E^c$ be the set of numbers in $[0,1]$ that \textit{do} contain at least one digit '4' in their decimal expansion. We will show $m(E^c) = 1$.
    
    \item \textbf{Step-by-Step Removal:}
    \begin{itemize}
        \item \textbf{1st digit:} Numbers with '4' in the first place are in $[0.4, 0.5)$. Length = $1/10$. There are 9 remaining intervals of length $1/10$ (like $[0.0, 0.1), [0.1, 0.2)$, etc.).
        \item \textbf{2nd digit:} In each of those 9 remaining intervals, there is a sub-interval of length $1/100$ that has a '4' in the second place (e.g., $[0.14, 0.15)$). Total length removed = $9 \times (1/100)$.
        \item \textbf{$k$-th digit:} At step $k$, we look at the intervals that haven't been "hit" yet. There are $9^{k-1}$ such intervals. From each, we remove a sub-interval of length $10^{-k}$.
    \end{itemize}

    \item \textbf{Sum the Measures:} The total measure removed is:
    \[ \sum_{k=1}^\infty \frac{9^{k-1}}{10^k} = \frac{1}{10} \sum_{k=1}^\infty \left(\frac{9}{10}\right)^{k-1} \]
    This is a geometric series with first term $a = 1/10$ and ratio $r = 9/10$.
    \[ \text{Total} = \frac{1/10}{1 - 9/10} = \frac{1/10}{1/10} = 1 \]

    \item \textbf{Conclusion:} Since the total length of $[0,1]$ is 1, and the measure of numbers \textit{with} a '4' is 1, the measure of $E$ must be $1 - 1 = 0$.
\end{enumerate}

\bigskip

\textbf{Intuition:} 
This set is essentially a "Fat Cantor Set" variant. Even though there are infinitely many numbers with no '4' (like $0.111\dots$ or $0.2323\dots$), they are so "thinly spread" that they take up no space at all. It's like a net with smaller and smaller holes—eventually, almost everything "falls through" the holes where the digit 4 appears.
\end{solution}

\begin{exercise}
\label{analysis:2012:10}
For $f \in L^p([0,1])$, show that $\lim_{\lambda\to\infty} \lambda^p m(\{|f|>\lambda\}) = 0$.
\end{exercise}

\begin{solution}
\textbf{Motivation:} This is a refinement of Markov's inequality. For $f \in L^p$, the "tail" of the integral must vanish.

1. $\int_0^1 |f|^p dm \geq \int_{\{|f|>\lambda\}} |f|^p dm \geq \lambda^p m(\{|f|>\lambda\})$.
2. Let $g = |f|^p$. Then $g \in L^1$. We want to show $\lambda^p m(\{g > \lambda^p\}) \to 0$.
3. Let $E_\lambda = \{|f| > \lambda\}$. As $\lambda \to \infty$, $E_\lambda \downarrow \emptyset$ a.e.
4. By the Dominated Convergence Theorem, $\int_{E_\lambda} |f|^p dm \to 0$ because $f \in L^p$ and the measure of $E_\lambda$ goes to $0$.
5. Since $\lambda^p m(E_\lambda) \leq \int_{E_\lambda} |f|^p dm$, the limit is 0.
\end{solution}

\begin{exercise}
\label{analysis:2012:11}
\begin{enumerate}
    \item Let $\mathfrak{F} = \{e_\nu\}$ be an orthonormal basis in an inner product space $H$. Let $E$ be the closed linear subspace spanned by $\mathfrak{F}$. For any $x \in H$, show that the following are equivalent:
    \begin{enumerate}
        \item $x \in E$;
        \item $\|x\|^2 = \sum_\nu |(x, e_\nu)|^2$;
        \item $x = \sum_\nu (x, e_\nu) e_\nu$.
    \end{enumerate}
    \item Let $H = L^2[0, 2\pi]$ with inner product $\langle f, g \rangle = \frac{1}{\pi} \int_0^{2\pi} f(x)g(x) dx$. Let
    \[ \mathfrak{F} = \left\{ \frac{1}{\sqrt{2}}, \cos x, \sin x, \dots, \cos nx, \sin nx, \dots \right\} \]
    be an orthonormal basis. Show that the closed linear subspace $E$ spanned by $\mathfrak{F}$ is $H$.
\end{enumerate}
\end{exercise}

\begin{solution}
\textbf{Motivation:} Completeness means the span is dense. We use the fact that continuous functions are dense in $L^2$ and trigonometric polynomials are dense in continuous periodic functions.

1. Continuous functions $C[0,2\pi]$ are dense in $L^2[0,2\pi]$.
2. For any $f \in C[0,2\pi]$, we can approximate it by periodic continuous functions by slightly modifying it near the boundaries.
3. By the Stone-Weierstrass Theorem, the span of $\{1, \cos nx, \sin nx\}$ is dense in $C_{per}[0,2\pi]$ under the sup-norm.
4. Since $\| \cdot \|_2 \leq \sqrt{2\pi} \| \cdot \|_\infty$, density in sup-norm implies density in $L^2$.
\end{solution}

\begin{exercise}
\label{analysis:2012:12}
Show that the Volterra operator $(Kf)(s) = \int_0^s f(t) dt$ is compact and has no eigenvalues.
\end{exercise}

\begin{solution}
\textbf{Motivation:} The Volterra operator is a classic example of a compact operator without eigenvalues (its spectrum is $\{0\}$, but it's not a nil-potent operator).

1. \textbf{Compactness:} Let $\|f_n\|_2 \leq 1$. Then $|Kf_n(s)| \leq 1$ and $|Kf_n(s_1) - Kf_n(s_2)| \leq \sqrt{|s_1-s_2|}$. By Arzel\`a-Ascoli, $\{Kf_n\}$ is precompact in $C[0,1]$, and thus in $L^2[0,1]$.
2. \textbf{Eigenvalues:} $Kf = \lambda f$.
   - If $\lambda=0$, $\int_0^s f = 0 \implies f=0$ a.e.
   - If $\lambda \neq 0$, $f(s) = \frac{1}{\lambda} \int_0^s f(t) dt$. This implies $f$ is continuous, then $C^1$.
   - Differentiating gives $f'(s) = \frac{1}{\lambda} f(s) \implies f(s) = C e^{s/\lambda}$.
   - At $s=0$, $f(0) = \frac{1}{\lambda} \int_0^0 f = 0$. Thus $C e^0 = 0 \implies C = 0 \implies f=0$.
\end{solution}
